\documentclass[11pt, a4paper]{article}
\usepackage[a4paper, top=2.5cm, bottom=2.5cm, left=2cm, right=2cm]{geometry}
\usepackage{iftex}
\ifPDFTeX
  \usepackage[utf8]{inputenc}
  \usepackage[T1]{fontenc}
  \usepackage[spanish, es-nodecimaldot, es-tabla]{babel}
\else
  \usepackage{fontspec}
  \usepackage[spanish, bidi=basic, provide=*]{babel}
  \babelprovide[import, onchar=ids fonts]{spanish}
  \babelprovide[import, onchar=ids fonts]{english}
  \babelfont{rm}{Noto Sans}
  \babelfont{tt}{Noto Sans Mono}
\fi
\usepackage{enumitem}
\setlist[itemize]{label=-}
\usepackage{graphicx}
\usepackage{float}
\usepackage{xcolor}
\usepackage{listings}
\usepackage{hyperref}

\definecolor{codebg}{RGB}{248, 249, 250}
\definecolor{codeframe}{RGB}{218, 220, 224}

\lstset{
  language=SQL,
  backgroundcolor=\color{codebg},
  frame=single,
  rulecolor=\color{codeframe},
  basicstyle=\ttfamily\small,
  breaklines=true,
  showstringspaces=false,
  tabsize=4,
  aboveskip=8pt,
  belowskip=8pt,
  literate={á}{{\'a}}1 {é}{{\'e}}1 {í}{{\'i}}1 {ó}{{\'o}}1 {ú}{{\'u}}1 {ñ}{{\~n}}1 {Ñ}{{\~N}}1
}

% Carga la imagen si está subida al proyecto (Overleaf); si no existe en el entorno, muestra el recuadro sin fallar
\newcommand{\mostrarcaptura}[2]{%
  \IfFileExists{#1}{%
    \includegraphics[width=0.92\textwidth]{#1}%
  }{%
    \framebox{\parbox{0.92\textwidth}{\centering
      \vspace{1.2cm}
      \textbf{Captura de MySQL Workbench: \texttt{#1}} \\
      \small\textit{#2}
      \vspace{1.2cm}
    }}%
  }%
}

\begin{document}

\begin{center}
    {\LARGE \textbf{Laboratorio 5: Entregable SQL (2/2)}}\\[0.2cm]
    {\Large Andes Bike}\\[0.3cm]
    {\normalsize ICN-292 Sistemas de Información para la Gestión $\cdot$ Segundo Semestre 2026}
\end{center}

\vspace{0.4cm}

\noindent
\begin{tabular}{@{}ll}
    \textbf{Estudiante:}  & Vicente Gabriel Soto Martínez \\
    \textbf{RUT:}         & 21.270.145-9 \\
    \textbf{Rol:}         & 202204651-1 \\
    \textbf{Profesores:}  & José Miguel González Paul, José Luis Sáez Tamayo, Patricio Cifuentes \\
    \textbf{Repositorio:} & \url{https://github.com/Vicente-Soto/ICN292-Lab5-Soto-Vicente}
\end{tabular}

\vspace{0.3cm}
\noindent\rule{\textwidth}{0.4pt}
\vspace{0.2cm}

\section*{1. Modelo entidad-relación}
El modelo se generó en MySQL Workbench a partir de las tablas de la base de datos.

\begin{figure}[H]
    \centering
    \IfFileExists{Diagrama.jpeg}{%
        \includegraphics[width=0.65\textwidth]{Diagrama.jpeg}%
    }{%
        \framebox{\parbox{0.65\textwidth}{\centering
            \vspace{1.2cm}
            \textbf{Captura de MySQL Workbench: \texttt{Diagrama.jpeg}} \\
            \small\textit{Modelo entidad-relación de la base de datos.}
            \vspace{1.2cm}
        }}%
    }
    \caption{Modelo entidad-relación de la base de datos.}
\end{figure}

\section*{2. Consultas}

\subsection*{Consulta 1}
\begin{lstlisting}
SELECT   v.nombre, COUNT(ve.id_venta) AS ventas
FROM     vendedores v
LEFT JOIN ventas ve ON v.id_vendedor = ve.id_vendedor
GROUP BY v.id_vendedor, v.nombre
ORDER BY ventas DESC, v.nombre ASC;
\end{lstlisting}

\begin{figure}[H]
    \centering
    \mostrarcaptura{Consulta1.png}{Resultado de la consulta 1 en MySQL Workbench.}
    \caption{Resultado de la consulta 1.}
\end{figure}

\subsection*{Consulta 2}
\begin{lstlisting}
SELECT   t.nombre,
         COUNT(ve.id_venta) AS ventas,
         SUM(ve.unidades)   AS unidades
FROM     tiendas    t
LEFT JOIN vendedores v  ON t.id_tienda   = v.id_tienda
LEFT JOIN ventas     ve ON v.id_vendedor = ve.id_vendedor
GROUP BY t.id_tienda, t.nombre
ORDER BY unidades DESC;
\end{lstlisting}

\begin{figure}[H]
    \centering
    \mostrarcaptura{Consulta 2.png}{Resultado de la consulta 2 en MySQL Workbench.}
    \caption{Resultado de la consulta 2.}
\end{figure}

\newpage

\subsection*{Consulta 3}
\begin{lstlisting}
SELECT t.nombre
FROM   tiendas    t
LEFT JOIN vendedores v ON t.id_tienda = v.id_tienda
WHERE  v.id_vendedor IS NULL;
\end{lstlisting}

\begin{figure}[H]
    \centering
    \mostrarcaptura{Consulta3.png}{Resultado de la consulta 3 en MySQL Workbench.}
    \caption{Resultado de la consulta 3.}
\end{figure}

\subsection*{Consulta 4}
\begin{lstlisting}
SELECT b.modelo
FROM   bicicletas b
LEFT JOIN ventas ve ON b.id_bicicleta = ve.id_bicicleta
                   AND ve.fecha BETWEEN '2025-06-01' AND '2025-06-30'
WHERE  ve.id_venta IS NULL;
\end{lstlisting}

\begin{figure}[H]
    \centering
    \mostrarcaptura{Consulta 4.png}{Resultado de la consulta 4 en MySQL Workbench.}
    \caption{Resultado de la consulta 4.}
\end{figure}

\noindent La condición del mes se puso en el ON del LEFT JOIN, porque si la poníamos en el WHERE no habría devuelto nada al pedir al mismo tiempo que tenga fecha en junio y que id\_venta sea NULL.

\newpage

\subsection*{Consulta 5}
\begin{lstlisting}
SELECT v.nombre, t.nombre AS tienda
FROM   vendedores v
LEFT JOIN tiendas t  ON v.id_tienda   = t.id_tienda
LEFT JOIN ventas  ve ON v.id_vendedor = ve.id_vendedor
WHERE  ve.id_venta IS NULL;
\end{lstlisting}

\begin{figure}[H]
    \centering
    \mostrarcaptura{Consulta 5.png}{Resultado de la consulta 5 en MySQL Workbench.}
    \caption{Resultado de la consulta 5.}
\end{figure}

\subsection*{Consulta 6}
\begin{lstlisting}
SELECT v1.nombre        AS vendedor_1,
       v1.fecha_ingreso AS ingreso_1,
       v2.nombre        AS vendedor_2,
       v2.fecha_ingreso AS ingreso_2,
       t.nombre         AS tienda
FROM   vendedores v1
JOIN   vendedores v2 ON v1.id_tienda     = v2.id_tienda
                    AND v1.fecha_ingreso < v2.fecha_ingreso
JOIN   tiendas    t  ON v1.id_tienda     = t.id_tienda;
\end{lstlisting}

\begin{figure}[H]
    \centering
    \mostrarcaptura{Consulta 6.png}{Resultado de la consulta 6 en MySQL Workbench.}
    \caption{Resultado de la consulta 6.}
\end{figure}

\newpage

\subsection*{Consulta 7}
\begin{lstlisting}
SELECT   t.nombre,
         SUM(ve.unidades * b.precio) AS monto_total
FROM     tiendas    t
JOIN     vendedores v  ON t.id_tienda    = v.id_tienda
JOIN     ventas     ve ON v.id_vendedor  = ve.id_vendedor
JOIN     bicicletas b  ON ve.id_bicicleta = b.id_bicicleta
WHERE    b.categoria IN ('Montaña', 'Ruta')
GROUP BY t.id_tienda, t.nombre
HAVING   SUM(ve.unidades * b.precio) > 2000000
ORDER BY monto_total DESC;
\end{lstlisting}

\begin{figure}[H]
    \centering
    \mostrarcaptura{Consulta 7.png}{Resultado de la consulta 7 en MySQL Workbench.}
    \caption{Resultado de la consulta 7.}
\end{figure}

\subsection*{Consulta 8}
\begin{lstlisting}
SELECT   modelo, precio
FROM     bicicletas
WHERE    precio > (SELECT AVG(precio)
                   FROM bicicletas)
ORDER BY precio DESC;
\end{lstlisting}

\begin{figure}[H]
    \centering
    \mostrarcaptura{Consulta 8.png}{Resultado de la consulta 8 en MySQL Workbench.}
    \caption{Resultado de la consulta 8.}
\end{figure}

\noindent El promedio es 657490 y lo obtuve corriendo aparte SELECT AVG(precio) FROM bicicletas (o sumando los 8 precios de la tabla bicicletas y dividiendo por 8).

\newpage

\subsection*{Consulta 9}
\begin{lstlisting}
SELECT   v.nombre, AVG(ve.unidades) AS promedio_unidades
FROM     vendedores v
JOIN     ventas     ve ON v.id_vendedor = ve.id_vendedor
GROUP BY v.id_vendedor, v.nombre
HAVING   AVG(ve.unidades) > (SELECT AVG(unidades)
                             FROM ventas)
ORDER BY promedio_unidades DESC, v.nombre ASC;
\end{lstlisting}

\begin{figure}[H]
    \centering
    \mostrarcaptura{Consulta 9.png}{Resultado de la consulta 9 en MySQL Workbench.}
    \caption{Resultado de la consulta 9.}
\end{figure}

\subsection*{Consulta 10}
\begin{lstlisting}
SELECT   v.nombre,
         SUM(ve.unidades * b.precio) AS monto_total
FROM     vendedores v
LEFT JOIN ventas     ve ON v.id_vendedor   = ve.id_vendedor
LEFT JOIN bicicletas b  ON ve.id_bicicleta = b.id_bicicleta
GROUP BY v.id_vendedor, v.nombre
ORDER BY monto_total DESC;
\end{lstlisting}

\begin{figure}[H]
    \centering
    \mostrarcaptura{Consulta 10.png}{Resultado de la consulta 10 en MySQL Workbench.}
    \caption{Resultado de la consulta 10.}
\end{figure}

\noindent No muestran cero porque como no tienen ventas el LEFT JOIN los deja en NULL y el SUM da NULL, para que salga 0 habría que poner IFNULL(SUM(ve.unidades * b.precio), 0).

\end{document}